% arXiv 投稿用 LaTeX 模板骨架
% 编译方式：pdflatex（正文使用英文；如需中文请改用 XeLaTeX + ctex）
% arXiv 投稿说明：https://info.arxiv.org/help/submit_tex.html
\documentclass[11pt]{article}

\usepackage[margin=1in]{geometry}
\usepackage{amsmath,amssymb}
\usepackage{graphicx}
\usepackage{booktabs}
\usepackage{caption}
\usepackage[colorlinks=true,linkcolor=blue,citecolor=blue,urlcolor=blue]{hyperref}

\title{Pathways to Building an Open-Source Aerospace Technology Sharing
Ecosystem and Deep-Space Simulation Platform Informed by the History of
Computer Development}

\author{Author Name\\
\small Institution, Country\\
\small \texttt{author@example.com}
}

\date{\today}

\begin{document}

\maketitle

\begin{abstract}
Space engineering is characterized by high technical barriers, heavy
investment, long development cycles and lengthy collaboration chains...
\end{abstract}

\noindent\textbf{Keywords:} history of computer development; open-source
ecosystem; aerospace technology sharing; deep-space simulation platform;
microservice architecture; DevSpaceOps; technology standardization; digital twin

% ---- 第 1 章 引言 ----
\section{Introduction}
% 1.1 研究背景与问题提出
% 1.2 研究意义
% 1.3 研究思路与方法
% 1.4 创新点

% ---- 第 2 章 文献综述与理论基础 ----
\section{Literature Review and Theoretical Foundations}
\subsection{The History of the Computer Industry}
\subsection{Open Source Communities and Open Innovation Theory}
\subsection{Open-Source Aerospace and Digital Twin Research}
\subsection{Gap Statement}

% ---- 第 3 章 计算机发展史的经验提炼 ----
\section{Lessons from the History of Computer Development}
\subsection{Open-Source Movement and Free Software}
\subsection{Standardization: From Proprietary to Open Standards}
\subsection{Modularity and Middleware: Decomposition of Complex Systems}
\subsection{Simulation and Virtualization: Low-Cost Iterative Validation}
\subsection{The Symbiosis of Open Source and Commerce}
\subsection{Common Patterns}
% 概念映射表（Table 1）
\begin{table}[htbp]
\centering
\caption{Concept mapping from computer history to open-source aerospace}
\begin{tabular}{@{}llll@{}}
\toprule
Prototype & Modern computing & Proposed aerospace form & Value \\
\midrule
OSI collaboration & Linux/Git/GitHub & OpenSpace ecosystem & Lower barriers \\
H/W--S/W decoupling & POSIX/HAL & Space OS & Cross-hardware reuse \\
Virtualization & QEMU/SITL & Deep-space digital twin & Sim-first validation \\
Modular services & Docker/microservices & Microservice spacecraft & Fault isolation \\
Delivery pipelines & CI/CD & DevSpaceOps & Automated reliability tests \\
Package ecosystems & npm/PyPI & Space Package Registry & Reuse, faster iteration \\
Open business loop & Red Hat/MySQL/Android & Open base + commercial services & Sustainable funding \\
Fault tolerance & BFT/zero-trust & Zero-trust space computing & Core safety baseline \\
\bottomrule
\end{tabular}
\end{table}

% ---- 第 4 章 开源航天生态总体架构 ----
\section{Architecture of the Open-Source Aerospace Ecosystem}
\subsection{Overall Architecture Blueprint}
\subsection{Open Collaboration Layer: OpenSpace R\&D Ecosystem}
\subsection{Standardization Layer: Open Interfaces and Protocols}
\subsection{Platform Layer: Universal Spacecraft Operating System (Space OS)}
\subsection{Reuse Layer: Space Package Registry}

% ---- 第 5 章 深空仿真平台 ----
\section{Building the Deep-Space Simulation Platform}
\subsection{Deep-Space Digital Twin Environment}
\subsection{DevSpaceOps: Cloud Simulation Driven High-Reliability R\&D}
\subsection{Shared Computing Resources and Federated Simulation}

% ---- 第 6 章 安全、信任与治理 ----
\section{Security, Trust and Governance}
\subsection{Zero-Trust Architecture and Code Security}
\subsection{Byzantine Fault Tolerance and Distributed Consensus}
\subsection{Community Governance and Accountability}
\subsection{The Tension Between Openness and Security}

% ---- 第 7 章 商业模式与可持续发展 ----
\section{Business Models and Sustainability}
\subsection{Dual Licensing and Commercial Open Source}
\subsection{Diversified Funding Sources}
\subsection{Industry--Academia--Research Collaboration and International Linkage}

% ---- 第 8 章 案例与可行性 ----
\section{Benchmark Cases and Feasibility Analysis}
\subsection{Benchmark Cases}
\subsection{Feasibility and Risk Analysis}

% ---- 第 9 章 结论与展望 ----
\section{Conclusion and Outlook}

\bibliographystyle{plain}
\bibliography{references}

\end{document}
